% ============================================================================
%  B/05-ba-tang-doc — file chương
%  Ngày tạo: 2026-09-06 | Sửa gần nhất: 2026-09-06 | Ôn gần nhất: chưa
%  Dependency: A/01, A/03. Mức độ: cốt lõi.
% ============================================================================
\documentclass[../../english.tex]{subfiles}
\begin{document}
\chapter{Ba tầng đọc và cách chọn tầng (Three Levels of Reading)}
\ptrtarget{B/05}\ptrtarget{B/05-ba-tang-doc}
\dependency{\ptr{A/01} cho ba tầng nghĩa; \ptr{A/03} cho kỹ năng tách câu --- cần khi xuống tới tầng đọc kỹ.}

\begin{tomtat}
Chương này bắt đầu bằng một câu hỏi mà người đọc tài liệu chuyên ngành hiếm khi tự hỏi: \emph{lần đọc này nhằm mục đích gì}. Trả lời câu đó trước sẽ quyết định bạn dùng tầng đọc nào trong ba tầng --- đọc thăm dò để quyết định có đáng đọc tiếp không, đọc kỹ để nắm nội dung và lập luận, và đọc phản biện để đánh giá. Chương đưa ra quy trình cho từng tầng, ngân sách thời gian gợi ý, một hệ thống ghi chú lề, và phép thử một câu để biết mình đã thật sự hiểu. Phần cuối nói về tầng thứ tư mà người làm nghiên cứu cần: đọc nhiều nguồn về cùng một chủ đề và so chúng với nhau. Nếu chỉ nhớ một điều: phần lớn thời gian đọc bị lãng phí không phải vì bạn đọc chậm mà vì bạn đọc sai tầng.
\end{tomtat}

\begin{nentang}
\kn{đọc thăm dò}{lướt có hệ thống qua tiêu đề, tóm tắt, đề mục, hình và kết luận để quyết định tài liệu này có đáng đọc kỹ không, và nếu có thì đọc phần nào.}
\kn{đọc kỹ}{đọc để nắm nội dung và bộ khung lập luận, có ghi chú, có dừng lại kiểm tra hiểu.}
\kn{đọc phản biện}{đọc để đánh giá: bằng chứng có đủ không, giả định nào chưa nói ra, kết luận có vượt quá dữ liệu không.}
\kn{đọc đối chiếu}{đọc nhiều nguồn về cùng một chủ đề và so chúng với nhau để dựng một bức tranh chung --- việc mà người làm nghiên cứu phải làm.}
\kn{ghi chú lề}{hệ thống ký hiệu ngắn viết ngay bên lề khi đọc, để lần sau quay lại không phải đọc lại toàn bộ.}
\kn{phép thử một câu}{gấp tài liệu lại và nói ý chính bằng một câu của mình; không làm được nghĩa là chưa hiểu, dù cảm giác nói khác.}
\end{nentang}

\begin{khoigoi}
\h{Vì sao ``đọc kỹ mọi thứ'' lại là một chiến lược tồi, kể cả khi bạn có nhiều thời gian?}
\h{Làm sao quyết định một bài báo có đáng đọc kỹ không, mà không phải đọc nó?}
\h{Đọc lướt và đọc thăm dò khác nhau thế nào --- và vì sao khác biệt đó quan trọng?}
\h{Bạn đọc xong một chương sách và thấy ``hiểu rồi''. Làm sao kiểm tra cảm giác đó có đúng không?}
\h{Khi cần nắm một chủ đề mới, đọc một nguồn thật kỹ hay đọc năm nguồn ở mức vừa phải?}
\end{khoigoi}

Hãy hình dung hai người cùng cần nắm một chủ đề mới trong hai tuần. Người thứ nhất tải về ba mươi bài báo, mở bài đầu tiên, và đọc từ đầu tới cuối, tra mọi từ lạ. Sau ba ngày anh ta xong bốn bài, mệt, và vẫn chưa biết bốn bài đó có phải bốn bài quan trọng nhất không. Người thứ hai dành hai giờ đầu tiên \emph{không đọc} bài nào cả --- cô ấy lướt qua cả ba mươi bài theo một quy trình cố định, loại ra hai mươi bài không liên quan, xếp mười bài còn lại theo thứ tự ưu tiên, và chỉ khi đó mới bắt đầu đọc.

Khác biệt giữa hai người không phải tốc độ đọc mà là việc người thứ hai coi ``đọc'' là nhiều hoạt động khác nhau chứ không phải một. Chương này làm rõ các hoạt động đó và cho tiêu chí chọn giữa chúng.

\section{Bốn câu hỏi trước khi đọc}

Trước khi mở một tài liệu, bốn câu hỏi này tốn ba mươi giây và quyết định phần còn lại.

\emph{Một --- tôi đọc cái này để làm gì?} Có bốn mục đích thường gặp, và chúng đòi bốn cách đọc: để quyết định có nên đọc tiếp không; để nắm ý chính; để làm theo được (cài đặt lại một phương pháp); hoặc để đánh giá và trích dẫn trong bài viết của mình.

\emph{Hai --- tôi đã biết gì về chủ đề này?} Nếu bạn đã biết nhiều, phần lớn tài liệu là thông tin cũ và bạn chỉ cần tìm phần mới. Nếu bạn hoàn toàn mới, hãy tìm một nguồn tổng quan trước --- đọc một bài báo chuyên sâu khi chưa có nền là cách chắc chắn nhất để mất thời gian.

\emph{Ba --- tôi cho nó bao nhiêu thời gian?} Đặt một ngân sách \emph{trước}. Không có ngân sách thì bạn sẽ đọc tới khi mệt chứ không tới khi đủ.

\emph{Bốn --- tôi cần giữ lại gì sau khi đọc?} Một câu tóm tắt? Một công thức? Một trích dẫn? Câu trả lời quyết định bạn ghi chú thế nào và ghi bao nhiêu.

\begin{cannho}
Sai lầm tốn kém nhất trong đọc tài liệu chuyên ngành không phải đọc chậm mà là \emph{đọc kỹ thứ không đáng}. Một giờ đọc kỹ một bài không liên quan tốn hơn nhiều so với mười phút thăm dò ba mươi bài. Vì vậy quy trình luôn bắt đầu bằng tầng một, kể cả khi bạn tin rằng tài liệu này chắc chắn quan trọng.
\end{cannho}

\section{Tầng một: đọc thăm dò}

Mục đích: trong năm tới mười phút, trả lời được ba câu --- tài liệu này nói về cái gì, nó khẳng định gì, và tôi có cần đọc kỹ không.

Quy trình sáu bước, theo đúng thứ tự này và \emph{không} đọc phần thân bài:

\begin{enumerate}
\item Đọc tiêu đề và tên tác giả. Nếu là bài báo, liếc năm xuất bản và nơi công bố --- chúng cho biết bối cảnh và mức độ nghiêm ngặt.
\item Đọc tóm tắt, chậm. Đây là phần dày thông tin nhất của cả bài.
\item Đọc mọi đề mục để dựng bộ khung.
\item Xem hình và bảng cùng chú thích --- trong bài kỹ thuật, hình thường mang kết quả chính, và chú thích được viết để đứng độc lập.
\item Đọc phần kết luận.
\item Liếc danh mục tài liệu tham khảo: bạn nhận ra bao nhiêu tên? Nếu không nhận ra ai, bạn đang ở ngoài vùng quen thuộc và nên tìm một bài tổng quan trước.
\end{enumerate}

Kết thúc sáu bước, hãy ra quyết định rõ ràng: bỏ, để dành, hay đọc kỹ. Việc \emph{bỏ} một tài liệu là một quyết định hợp lệ và cần được luyện --- rất nhiều người đọc kém hiệu quả chỉ vì không cho phép mình bỏ.

Chú ý phân biệt \emph{đọc thăm dò} với \emph{đọc lướt}. Đọc lướt là đọc nhanh toàn văn, mắt chạy qua mọi dòng; nó vừa mệt vừa cho hiểu biết mờ. Đọc thăm dò là đọc \emph{đủ chậm} nhưng chỉ đọc những phần mang thông tin cấu trúc. Đây là câu trả lời cho câu hỏi mở đầu thứ ba, và khác biệt này thay đổi hẳn hiệu quả.

\section{Tầng hai: đọc kỹ}

Mục đích: nắm nội dung và bộ khung lập luận đủ để giải thích lại cho người khác.

Ba nguyên tắc.

\emph{Nguyên tắc một --- đọc theo đơn vị đoạn, không theo câu.} Mỗi đoạn trong văn học thuật thường có một ý chính, và ý đó hay nằm ở câu đầu hoặc câu cuối. Sau mỗi đoạn, dừng nửa giây và tự hỏi ``đoạn này làm gì trong lập luận'' --- nêu vấn đề, đưa bằng chứng, phản bác, hay chuyển ý. Chương \ptr{B/06} cho bộ từ vựng để trả lời câu đó nhanh.

\emph{Nguyên tắc hai --- ghi chú lề bằng ký hiệu, không bằng câu.} Một hệ thống tối giản dùng được ngay: \textbf{?} chỗ chưa hiểu; \textbf{!} chỗ bất ngờ hoặc quan trọng; \textbf{$\to$} chỗ nối với thứ bạn đã biết; \textbf{C} chỗ tác giả nêu khẳng định chính; \textbf{E} chỗ đưa bằng chứng; \textbf{L} chỗ thừa nhận hạn chế. Ghi chú bằng ký hiệu nhanh hơn viết câu nên không làm đứt mạch, và khi quay lại bạn tìm được ngay phần cần.

\emph{Nguyên tắc ba --- dừng lại ở chỗ khó, nhưng có giới hạn.} Khi gặp một đoạn không hiểu, hãy thử ba việc theo thứ tự: đọc lại chậm và tách câu theo \ptr{A/03}; đọc tiếp một đoạn nữa vì tác giả thường giải thích ngay sau; và đánh dấu \textbf{?} rồi đi tiếp. Đừng dừng quá hai phút ở một đoạn trong lần đọc đầu --- rất nhiều chỗ khó tự sáng ra khi bạn có bức tranh toàn cục.

Cuối lần đọc kỹ, làm \emph{phép thử một câu}: gấp tài liệu và nói ý chính bằng một câu của mình. Đây là câu trả lời cho câu hỏi mở đầu thứ tư. Nếu bạn chỉ lặp lại được các từ khoá của tác giả mà không diễn đạt lại được, bạn mới nhận ra chứ chưa hiểu --- đúng phân biệt ở \ptr{A/01}.

\section{Tầng ba: đọc phản biện}

Mục đích: quyết định tin bao nhiêu, và dùng được kết quả này vào việc gì.

Sáu câu hỏi, dùng cho mọi tài liệu có đưa ra khẳng định thực nghiệm:

\begin{enumerate}
\item Khẳng định chính là gì, phát biểu bằng một câu?
\item Bằng chứng nào được đưa ra cho nó?
\item Bằng chứng đó có đủ để suy ra khẳng định không --- hay có một bước nhảy ở giữa?
\item Giả định nào chưa được nói ra? (Điều gì phải đúng thì lập luận mới chạy?)
\item Kết luận có vượt quá phạm vi dữ liệu không? Tìm các cụm giới hạn phạm vi ở \ptr{A/01}.
\item Ai có thể phản bác điều này, và tác giả đã trả lời chưa?
\end{enumerate}

Bốn dấu hiệu đáng ngờ cần để ý: khẳng định mạnh mà không có số liệu; so sánh với các phương pháp cũ nhưng không nói rõ điều kiện so sánh; kết quả chỉ báo cáo ở trường hợp tốt nhất; và phần hạn chế viết chung chung tới mức không nói gì. Chương \ptr{B/08} cho công cụ tinh hơn: đọc mức cam kết qua chính cách tác giả chọn từ.

\section{Tầng bốn: đọc đối chiếu nhiều nguồn}

Đây là tầng mà người làm nghiên cứu dùng nhiều nhất nhưng ít được dạy nhất, và nó trả lời câu hỏi mở đầu thứ năm.

Khi cần nắm một chủ đề mới, đọc \emph{năm nguồn ở mức vừa} gần như luôn tốt hơn đọc \emph{một nguồn thật kỹ}. Ba lý do. Thứ nhất, mỗi tác giả có thiên hướng riêng và bạn chỉ thấy được điều đó khi có cái để so. Thứ hai, những gì \emph{lặp lại} ở nhiều nguồn chính là phần đã ổn định của lĩnh vực, còn những gì chỉ một nguồn nói là phần còn tranh cãi --- một phân biệt rất quan trọng mà đọc một nguồn không cho bạn. Thứ ba, từ vựng và cách diễn đạt lặp lại, nên bạn học ngôn ngữ nhanh hơn nhiều --- đúng chiến lược đọc hẹp ở \ptr{A/02}.

Quy trình đọc đối chiếu: chọn một câu hỏi cụ thể; thu thập năm tới mười nguồn; đọc thăm dò tất cả để loại và xếp thứ tự; đọc kỹ ba tới bốn nguồn tốt nhất; rồi lập một bảng so sánh với các cột là tiêu chí bạn quan tâm và các hàng là nguồn. Chính cái bảng đó --- chứ không phải các ghi chú riêng lẻ --- là thứ cho bạn bức tranh chung.

\begin{cambay}
Cạm bẫy khi đọc nhiều nguồn: tưởng rằng cùng một thuật ngữ ở hai bài có cùng nghĩa. Trong các lĩnh vực trẻ, cùng một từ thường được định nghĩa khác nhau, và cùng một khái niệm được gọi bằng nhiều tên. Khi lập bảng so sánh, hãy thêm một cột ghi rõ \emph{mỗi tác giả định nghĩa thuật ngữ then chốt thế nào} --- chương \ptr{B/10} nói kỹ về hiện tượng này và cách xử lý.
\end{cambay}

\section{Gốc rễ: từ đọc như một kỹ năng tới đọc như nhiều kỹ năng}

Ý tưởng rằng đọc không phải một hoạt động đồng nhất được trình bày có hệ thống lần đầu trong một quyển sách xuất bản năm 1940 và sửa lại năm 1972: Adler và Van Doren phân biệt bốn mức đọc, trong đó ba mức trên gần trùng với ba tầng ở chương này, còn mức thứ tư --- đọc nhiều sách về cùng chủ đề để tự dựng câu trả lời --- chính là tầng đối chiếu ở mục 5\cite{adler1972}. Điểm đáng chú ý về mặt lịch sử: quyển sách ra đời trong bối cảnh giáo dục đại chúng mở rộng nhanh và người ta nhận ra rằng biết chữ không đồng nghĩa với biết đọc một cuốn sách khó.

Cùng thời kỳ, trong giáo dục Mỹ xuất hiện các phương pháp đọc có quy trình --- nổi tiếng nhất là sơ đồ khảo sát, đặt câu hỏi, đọc, kể lại, ôn tập. Điểm chung của chúng với chương này là nguyên tắc \emph{khảo sát trước khi đọc}: dựng bộ khung trong đầu rồi mới nạp chi tiết vào khung đó.

Về sau, khi khối lượng tài liệu khoa học tăng vọt, các cộng đồng nghiên cứu tự phát triển những quy trình riêng cho việc đọc bài báo --- nổi tiếng nhất là phương pháp ba lượt được trình bày trong một bài viết ngắn năm 2007 và trở thành tài liệu được truyền tay trong giới nghiên cứu máy tính\cite{keshav2007}. Chương \ptr{B/07} dành trọn cho nó.

Điểm chung của cả ba dòng: chúng đều xuất phát từ một tình huống thừa thãi thông tin, và đều đưa ra cùng một lời giải --- \emph{quyết định trước khi đọc}, thay vì đọc rồi mới quyết định.

\begin{gocre}
\gr{1940, sửa 1972 --- Adler \& Van Doren}{Giáo dục mở rộng, nhiều người biết chữ nhưng không đọc nổi sách khó.}{Bốn mức đọc; ý tưởng rằng đọc là nhiều kỹ năng khác nhau, chọn theo mục đích.}
\gr{Giữa thế kỷ 20 --- các phương pháp đọc có quy trình}{Sinh viên đọc giáo trình mà không nhớ được gì.}{Nguyên tắc khảo sát trước khi đọc: dựng khung rồi mới nạp chi tiết --- nền của tầng một ở mục 2.}
\gr{2007 --- Keshav}{Nghiên cứu sinh phải đọc hàng chục bài báo mỗi tuần.}{Phương pháp ba lượt với ngân sách thời gian rõ ràng cho mỗi lượt (\ptr{B/07}).}
\gr{Thập niên 2010--nay --- bùng nổ số lượng công bố}{Số bài báo tăng nhanh hơn khả năng đọc của bất kỳ ai.}{Kỹ năng \emph{loại bỏ} trở nên quan trọng ngang kỹ năng đọc; đọc thăm dò thành bước bắt buộc.}
\end{gocre}

\section{Liên hệ ngang}

\begin{lienhe}
\lh{Đọc mã nguồn}{Xem kiến trúc và tên hàm trước, rồi mới đọc chi tiết}{Cùng nguyên tắc dựng khung trước. Lập trình viên có kinh nghiệm không đọc mã từ dòng một --- họ tìm điểm vào, xem cấu trúc thư mục, rồi mới đi sâu vào phần cần.}
\lh{Y khoa}{Phân loại bệnh nhân theo mức ưu tiên}{Đọc thăm dò chính là phân loại: quyết định nhanh cái gì cần xử lý kỹ, cái gì có thể chờ. Cùng một bài học: dành thời gian đều cho mọi thứ là cách chắc chắn để làm hỏng việc quan trọng.}
\lh{Nghiên cứu tổng quan}{Tổng quan hệ thống trong y học}{Tầng đọc đối chiếu ở mức thể chế hoá: tiêu chí chọn nguồn được ghi rõ trước, và kết quả là bảng so sánh có cấu trúc --- đúng thứ mục 5 khuyến nghị làm ở quy mô cá nhân.}
\lh{Tình báo và phân tích}{Đọc để trả lời một câu hỏi cụ thể, không đọc để biết chung}{Cùng nguyên tắc ``xác định câu hỏi trước''. Câu hỏi càng cụ thể thì việc loại bỏ tài liệu càng dễ và càng nhanh.}
\lh{Toán học}{Đọc chứng minh: nắm sơ đồ trước, kiểm chi tiết sau}{Người đọc chứng minh có kinh nghiệm đọc theo hai lượt --- lượt đầu nắm ý tưởng, lượt sau kiểm từng bước. Đây là ba tầng của chương này áp cho một thể loại rất hẹp.}
\lh{Quản lý thông tin cá nhân}{Ghi chú liên kết, thẻ ghi chú}{Ghi chú lề ở mục 3 là bước đầu; hệ thống ghi chú liên kết\cite{ahrens2017} là bước sau, biến các ghi chú rời thành mạng dùng lại được (\ptr{B/07}).}
\end{lienhe}

\begin{traloi}
\tl{Vì thời gian đọc là tài nguyên hữu hạn còn tài liệu thì không. Đọc kỹ mọi thứ nghĩa là dành thời gian như nhau cho tài liệu quan trọng và tài liệu không liên quan, và vì phần lớn tài liệu \emph{không} liên quan tới câu hỏi của bạn nên phần lớn thời gian bị lãng phí. Ngoài ra, đọc kỹ khi chưa có bức tranh tổng thể còn kém hiệu quả về mặt hiểu: nhiều chi tiết chỉ có nghĩa khi bạn đã biết chúng nằm ở đâu trong bộ khung.}
\tl{Bằng quy trình sáu bước ở mục 2, chỉ đọc phần mang thông tin cấu trúc: tiêu đề, tóm tắt, đề mục, hình và chú thích, kết luận, và danh mục tham khảo. Sáu bước này mất năm tới mười phút và trả lời được ba câu: bài nói về gì, khẳng định gì, và có liên quan tới câu hỏi của bạn không. Điều quan trọng là kết thúc bằng một \emph{quyết định rõ ràng} --- bỏ, để dành, hay đọc kỹ --- chứ không để mở.}
\tl{Đọc lướt là chạy mắt nhanh qua \emph{toàn văn}; nó tốn sức, cho hiểu biết mờ, và bạn không biết mình đã bỏ sót gì. Đọc thăm dò là đọc \emph{chậm và kỹ} nhưng chỉ ở những phần mang thông tin cấu trúc. Khác biệt quan trọng vì nó đổi bản chất công việc: đọc lướt là cố nắm nội dung với ít thời gian; đọc thăm dò là thu thập đủ thông tin để \emph{ra một quyết định}. Mục tiêu khác nhau nên cách làm khác nhau.}
\tl{Bằng phép thử một câu: gấp tài liệu lại và phát biểu ý chính bằng một câu của \emph{chính bạn}, không dùng lại các từ khoá của tác giả. Nếu chỉ lặp lại được thuật ngữ gốc thì đó là nhận ra chứ không phải hiểu --- đúng phân biệt ở \ptr{A/01} và \ptr{A/02}. Phép thử mạnh hơn nữa: giải thích cho một người không cùng chuyên môn, hoặc viết ba câu tóm tắt rồi so lại với bản gốc sau một ngày.}
\tl{Đọc năm nguồn ở mức vừa, gần như luôn. Ba lý do: bạn chỉ nhận ra thiên hướng của một tác giả khi có nguồn khác để so; những gì lặp lại ở nhiều nguồn là phần đã ổn định còn những gì chỉ một nguồn nói là phần còn tranh cãi; và từ vựng lặp lại giúp bạn học ngôn ngữ nhanh hơn hẳn. Ngoại lệ: khi có một tài liệu chuẩn mực của lĩnh vực --- một giáo trình kinh điển --- thì đọc kỹ nó trước rồi mới đối chiếu là hợp lý hơn.}
\end{traloi}

\begin{dongian}
Nhiều người coi đọc là một việc duy nhất: mở tài liệu ra và đọc từ đầu tới cuối. Đó là lý do họ mệt và không nhớ được bao nhiêu.

Thực ra có ít nhất ba việc khác nhau, và chúng khác nhau ngay từ mục đích. Việc thứ nhất là \emph{xem thử}: giống như cầm quyển sách trong hiệu sách, đọc bìa sau, lật mục lục, xem vài trang, rồi quyết định có mua không. Với bài báo, việc này mất năm phút và gồm: đọc tóm tắt, xem đề mục, nhìn hình, đọc kết luận. Kết thúc là một quyết định dứt khoát --- bỏ, để dành, hay đọc kỹ. Cho phép mình \emph{bỏ} là điều khó nhất và cũng cần nhất.

Việc thứ hai là \emph{đọc kỹ} để hiểu. Ở đây nên đọc theo đoạn chứ không theo câu, và sau mỗi đoạn tự hỏi ``đoạn này để làm gì'' --- nêu vấn đề, đưa bằng chứng, hay phản bác. Gặp chỗ khó thì đừng đứng lại quá hai phút; đánh dấu và đi tiếp, vì rất nhiều chỗ khó tự sáng ra khi bạn đã thấy toàn cảnh.

Việc thứ ba là \emph{đọc để đánh giá}: tác giả khẳng định gì, đưa bằng chứng gì, và giữa hai cái đó có bước nhảy nào không.

Cuối cùng, khi cần nắm một chủ đề mới, đừng đọc một tài liệu thật kỹ. Hãy đọc năm tài liệu ở mức vừa rồi lập một bảng so sánh. Cái lặp lại ở nhiều nguồn là phần đã chắc chắn; cái chỉ một nguồn nói là phần còn đang tranh cãi. Bạn không thể thấy sự phân biệt đó nếu chỉ đọc một nguồn, dù đọc kỹ tới đâu.
\motcau{Trước khi đọc, hãy quyết định lần đọc này để làm gì --- và đọc kỹ thứ không đáng còn tốn hơn đọc chậm.}
\chuadongian{Không có công thức chung cho việc ``bao nhiêu là đủ kỹ''; nó phụ thuộc mục đích, và chỉ kinh nghiệm mới cho cảm giác đúng.}
\end{dongian}

\begin{onbai}
\ar[3]{Bốn câu hỏi trước khi đọc}{Đọc để làm gì; tôi đã biết gì; cho nó bao nhiêu thời gian; cần giữ lại gì sau khi đọc.}
\ar[3]{Sáu bước đọc thăm dò}{Tiêu đề và nguồn \(\rightarrow\) tóm tắt \(\rightarrow\) đề mục \(\rightarrow\) hình và chú thích \(\rightarrow\) kết luận \(\rightarrow\) danh mục tham khảo; kết thúc bằng quyết định bỏ / để dành / đọc kỹ.}
\ar[3]{Đọc thăm dò khác đọc lướt}{Lướt là chạy mắt qua toàn văn (mờ và mệt); thăm dò là đọc kỹ nhưng chỉ ở phần mang thông tin cấu trúc, nhằm ra một quyết định.}
\ar[3]{Ba nguyên tắc đọc kỹ}{Đọc theo đoạn và hỏi ``đoạn này làm gì''; ghi chú lề bằng ký hiệu; dừng ở chỗ khó nhưng không quá hai phút.}
\ar[3]{Phép thử một câu}{Gấp tài liệu, nói ý chính bằng câu của mình, không dùng từ khoá của tác giả. Không làm được là chưa hiểu.}
\ar[3]{Sáu câu hỏi đọc phản biện}{Khẳng định chính; bằng chứng; có bước nhảy không; giả định chưa nói; kết luận có vượt phạm vi không; ai phản bác được và tác giả đã trả lời chưa.}
\ar[2]{Hệ ký hiệu ghi chú lề}{\textbf{?} chưa hiểu, \textbf{!} bất ngờ, \(\rightarrow\) nối với cái đã biết, \textbf{C} khẳng định, \textbf{E} bằng chứng, \textbf{L} hạn chế.}
\ar[2]{Đọc đối chiếu nhiều nguồn}{Năm nguồn mức vừa hơn một nguồn thật kỹ: thấy được thiên hướng, phân biệt phần ổn định với phần tranh cãi, và học ngôn ngữ nhanh hơn.}
\ar[2]{Quy trình đọc đối chiếu}{Câu hỏi cụ thể \(\rightarrow\) thu thập 5--10 nguồn \(\rightarrow\) thăm dò tất cả để xếp hạng \(\rightarrow\) đọc kỹ 3--4 \(\rightarrow\) lập bảng so sánh có cột định nghĩa thuật ngữ.}
\ar[2]{Bốn dấu hiệu đáng ngờ}{Khẳng định mạnh không số liệu; so sánh không rõ điều kiện; chỉ báo cáo trường hợp tốt nhất; phần hạn chế viết chung chung.}
\ar[1]{Vì sao cho phép mình bỏ tài liệu}{Kỹ năng loại bỏ quan trọng ngang kỹ năng đọc khi số tài liệu vượt xa khả năng đọc của bất kỳ ai.}
\end{onbai}

\begin{hoidap}
\qa{Tôi luôn thấy có lỗi khi bỏ dở một tài liệu. Làm sao vượt qua?}{Hãy đổi cách nghĩ về đơn vị công việc: mục tiêu không phải ``đọc xong tài liệu này'' mà ``trả lời câu hỏi của tôi''. Với khung đó, bỏ một tài liệu không liên quan là một quyết định \emph{thành công}, không phải thất bại --- bạn vừa tiết kiệm một giờ cho tài liệu đáng đọc hơn. Một mẹo thực dụng: ghi lại lý do bỏ vào một dòng ngắn. Việc ghi ra biến quyết định từ cảm tính thành có cơ sở, và nó cũng giúp bạn không mở lại cùng tài liệu đó ba tuần sau.}
\qa{Đọc thăm dò có bỏ sót tài liệu quan trọng không?}{Có, nhưng ít hơn bạn lo, và cái giá nhỏ hơn nhiều so với lợi ích. Tài liệu thực sự quan trọng thường được nhắc lại trong các nguồn khác, nên bạn sẽ gặp lại nó. Để giảm rủi ro: khi thăm dò, hãy đọc kỹ phần tóm tắt và \emph{đặc biệt là hình} --- nếu một hình cho thấy kết quả liên quan trực tiếp tới câu hỏi của bạn thì giữ bài đó lại dù tóm tắt viết không hấp dẫn.}
\qa{Ghi chú thế nào để sau ba tháng còn dùng được?}{Nguyên tắc: ghi chú phải đứng độc lập, hiểu được mà không cần mở lại tài liệu gốc. Ba thành phần tối thiểu cho mỗi tài liệu: một câu về khẳng định chính; hai tới ba dòng về bằng chứng và điều kiện; và một dòng về ``điều này liên quan gì tới việc tôi đang làm''. Dòng thứ ba là dòng khiến ghi chú sống được lâu, và nó cũng là cầu nối giữa các tài liệu khác nhau\cite{ahrens2017}.}
\qa{Tôi đọc tiếng Anh chậm hơn tiếng Việt nhiều. Có nên luyện đọc nhanh không?}{Các kỹ thuật đọc nhanh cực đoan --- bỏ đọc thầm, mở rộng tầm mắt --- có bằng chứng khá yếu, và với tài liệu chuyên ngành thì chúng đánh đổi hiểu lấy tốc độ, tức đánh đổi sai thứ. Tốc độ đọc tiếng Anh tăng chủ yếu nhờ hai thứ khác: độ phủ từ vựng cao hơn (\ptr{A/02}) và khả năng nhận ra cấu trúc câu tự động (\ptr{A/03}). Cả hai đều là kết quả của việc đọc nhiều trong cùng một lĩnh vực, không phải của kỹ thuật quét mắt.}
\qa{Khi nào nên đọc bản dịch tiếng Việt thay vì bản gốc?}{Khi mục tiêu là \emph{nắm khái niệm} mới hoàn toàn và bạn đang mất quá nhiều sức cho ngôn ngữ --- lúc đó đọc tiếng Việt trước rồi quay lại bản gốc là chiến lược tốt, vì lần đọc thứ hai bạn đã có khung khái niệm. Không nên dùng bản dịch khi mục tiêu là \emph{trích dẫn hoặc đánh giá chính xác}, vì bản dịch thường làm mất các sắc thái cam kết mà \ptr{B/08} sẽ nói tới --- và đó chính là chỗ chứa nhiều thông tin nhất.}
\qa{Đọc trên giấy hay trên màn hình?}{Với đọc thăm dò thì màn hình tiện hơn vì bạn cần nhảy qua lại. Với đọc kỹ một tài liệu quan trọng, nhiều người thấy giấy hoặc máy đọc chuyên dụng giúp tập trung hơn và ghi chú lề tự nhiên hơn --- và bằng chứng nghiên cứu cũng nghiêng nhẹ về phía đó cho việc hiểu văn bản dài. Điều chắc chắn hơn cả lựa chọn phương tiện: hãy tắt thông báo, vì đọc sâu bị phá hỏng bởi ngắt quãng nhiều hơn bởi định dạng.}
\qa{Làm sao biết mình đã đọc đủ về một chủ đề?}{Ba dấu hiệu thực dụng. Một, bạn bắt đầu gặp lại cùng các tên và cùng các ý ở nguồn mới --- dấu hiệu bão hoà. Hai, bạn dự đoán được nội dung phần tiếp theo trước khi đọc. Ba, và quan trọng nhất, bạn viết được một đoạn tổng hợp về chủ đề đó bằng lời của mình, nêu được cả những điểm còn tranh cãi. Nếu chưa viết được thì chưa đủ --- và bản thân việc viết sẽ chỉ ra chỗ còn thiếu, đúng tinh thần phần C.}
\end{hoidap}

\begin{baitap}
\bt[1]{Thăm dò năm bài báo trong 30 phút và ra quyết định cho từng bài}{Tài liệu ngành}{Ghi một dòng lý do cho mỗi quyết định bỏ / để dành / đọc kỹ.}
\bt[1]{Đọc kỹ một bài đã chọn và làm phép thử một câu}{Tài liệu ngành}{Viết câu tóm tắt ra giấy \emph{trước} khi mở lại bài để kiểm.}
\bt[2]{Dùng hệ ký hiệu ghi chú lề cho một bài dài}{Tài liệu ngành}{Sau đó chỉ đọc lại các chỗ đánh dấu và xem có dựng lại được lập luận không.}
\bt[2]{Áp sáu câu hỏi đọc phản biện cho một bài bạn thấy thuyết phục}{Tài liệu ngành}{Mục tiêu là tìm ra ít nhất một giả định chưa được nói ra.}
\bt[2]{Viết ghi chú ba thành phần cho ba tài liệu}{Bài tập ghi chú}{Khẳng định chính; bằng chứng và điều kiện; liên quan gì tới việc của bạn.}
\bt[3]{Đọc đối chiếu: chọn một câu hỏi hẹp, đọc năm nguồn, lập bảng so sánh}{Tài liệu ngành}{Cột bắt buộc: định nghĩa thuật ngữ then chốt của mỗi tác giả.}
\bt[3]{Sau một tuần, đọc lại ghi chú cũ và tự chấm chúng còn dùng được không}{Ghi chú của bạn}{Ghi chú nào phải mở lại bản gốc mới hiểu thì cần viết lại theo nguyên tắc ở mục hỏi đáp.}
\end{baitap}

\section*{Kết chương và đọc tiếp}

Chương này chuyển việc đọc từ một hoạt động đồng nhất thành bốn hoạt động có mục đích khác nhau, mỗi cái có quy trình riêng. Ba điều đáng mang theo: bốn câu hỏi trước khi đọc; sáu bước thăm dò kết thúc bằng một quyết định dứt khoát; và phép thử một câu để phân biệt hiểu thật với cảm giác quen.

Ở tầng đọc kỹ, chương này đã nhắc tới việc ``hỏi mỗi đoạn làm gì trong lập luận'' nhưng chưa cho công cụ để trả lời nhanh. Đó là nội dung của \ptr{B/06}: bộ khung lập luận học thuật, các từ nối báo hiệu quan hệ giữa các ý, và cách vẽ lại lập luận của một bài thành một sơ đồ mà bạn kiểm tra được.
\end{document}
